\section{Problem 2}

\subsection{Post Selected Pauli Projection}

\subsubsection{Answer}


Note we have the following identity

\begin{align*}
	\tikzfig{q-2-1.1} 
\end{align*}

So we have 
\begin{equation}\label{q2-1}
	W=\tikzfig{q-2-1.2}
\end{equation}

where $\vec{Q} = Z \otimes \vec{P}$ and is the desired answer.


\subsubsection{Analysis}

Let us denote the whole gadget as $W$ and analyse how does \eqref{q2-1} work.
We know

$$
\tikzfig{q-2-1.3} = |0\rangle + e^{\frac{\pi}{4}i} |1\rangle
$$

Let $v$ denote eigenvector of $\vec{P}$ with eigenvalue $1$, $w$ the eigenvector with eigenvalue $-1$; then $\vec{Q}$ has eigenvector $|0\rangle \otimes v$ and $|1\rangle \otimes w$ with eigenvalue $1$, and eigenvector $|1\rangle \otimes v$ and $|0\rangle \otimes w$ with eigenvalue $-1$. 
(This is because we added a $Z$ in the first qubit, which has eigenvalue $1$ for $|0\rangle$ and $-1$ for $|1\rangle$.)

The measurement $\Pi^0 \vec{Q}$ projects the state to the eigenspace of $\vec{Q}$ with eigenvalue $1$. 
while $\Pi^1 \vec{Q}$ projects the state to the eigenspace with eigenvalue $-1$.
Thus
\begin{equation}\label{q2-1-cal}
\begin{split}
	\Pi^0 \vec{Q} \left( \tikzfig{q-2-1.3} \otimes v \right) &= |0\rangle \otimes v \\
	\Pi^0 \vec{Q} \left( \tikzfig{q-2-1.3} \otimes w \right) &= e^{\frac{\pi}{4}i}|1\rangle \otimes w \\
	\Pi^1 \vec{Q} \left( \tikzfig{q-2-1.3} \otimes v \right) &= e^{\frac{\pi}{4}i}|1\rangle \otimes v \\
	\Pi^1 \vec{Q} \left( \tikzfig{q-2-1.3} \otimes w \right) &= |0\rangle \otimes v 
\end{split}
\end{equation}
This means
$$
W(v) = v, W(w) = e^{\frac{\pi}{4}i} w
$$

Generalising our discussion above gives the following lemma.

\begin{lemma}\label{lm:general-phase-pauli}
	Let $M$ denote the following map 
	$$
	M = \tikzfig{general-phase-pauli}
	$$
	Where $\vec{P}$ is a stabiliser Pauli. 
	Assuming $v$ is an eigenvector of $\vec{P}$ with eigenvalue $1$, and $w$ is an eigenvector of $\vec{P}$ with eigenvalue $-1$, we have
$$
M(v) = v, M(w) = e^{\alpha i} w
$$
\end{lemma}


